\begin{lemma}
\label{editorial-source-lemma-geometrically-irreducible-add-transcendental}
Let $K/L/M$ be a tower of fields with $L/M$ geometrically irreducible.
Let $x \in K$ be transcendental over $L$. Then $L(x)/M(x)$ is geometrically
irreducible.
\end{lemma}

\begin{proof}
Keep the original inclusions
$M\subseteq L\subseteq K$ and
the original $x\in K$.
Transcendence over $L$ makes
evaluation at $x$ injective on
$L[t]$, and hence on its subring
$M[t]$. In particular no original
nonzero denominator vanishes.
The maps
$$
\lambda:L(t)\longrightarrow L(x),\quad
\frac{P(t)}{Q(t)}\longmapsto\frac{P(x)}{Q(x)},
\qquad
\mu:M(t)\longrightarrow M(x)
$$
with the same evaluation formula
are injective field maps. They
are surjective because their
images are fields containing the
original coefficient fields and
$x$, respectively. Their inverses
replace $x$ by $t$ in a rational
expression; injectivity proves
independence of that expression.
The square with the original
inclusions commutes because the
formulas agree on every $M$
coefficient, on $t$, and on
every inverse denominator.

For any original field $E/M(x)$,
use its given embedding composed
with $\mu$ to view it over $M(t)$.
The map
$$
L(t)\otimes_{M(t)}E\longrightarrow L(x)\otimes_{M(x)}E,
\qquad a\otimes b\longmapsto\lambda(a)\otimes b
$$
is well-defined by that commuting
square, with inverse
$c\otimes b\mapsto\lambda^{-1}(c)\otimes b$.
Lemma
\ref{editorial-source-lemma-geometrically-irreducible-base-change-transcendental}
makes the left-hand spectrum
irreducible from the hypothesis
on $L/M$. This proves the assertion
for every given $E/M(x)$.
Conversely the same tensor maps,
using $\mu^{-1}$ to transport a
field over $M(t)$, show that
geometric irreducibility of
$L(x)/M(x)$ implies that of
$L(t)/M(t)$. The converse in the
preceding lemma then gives that
of $L/M$. Thus, under the stated
transcendence hypothesis, the
original conclusion is in fact
equivalent to geometric
irreducibility of $L/M$.
\end{proof}